\chapter{Lezione 19 - 03/12}

\section{Posizione Reciproca di Sottospazi Affini}

\Definizione{Sottospazi Paralleli}{
    Siano \(\mathbb{H}, \mathbb{H}'\) due sottospazi affini di uno spazio \(\mathcal{E}\).
    Si dice che \(\mathbb{H}\) e \(\mathbb{H}'\) sono \textbf{paralleli} (\(\mathbb{H} \parallel \mathbb{H}'\)) se le loro giaciture sono una contenuta nell'altra:
    \[
    \vec{\mathbb{H}} \subseteq \vec{\mathbb{H}'} \quad \text{oppure} \quad \vec{\mathbb{H}'} \subseteq \vec{\mathbb{H}}.
    \]
}

\Definizione{Sottospazi Sghembi}{
    Siano \(\mathbb{H}, \mathbb{H}'\) due sottospazi affini.
    \begin{itemize}
        \item \(\mathbb{H}\) e \(\mathbb{H}'\) si dicono \textbf{sghembi} se non hanno punti in comune e non sono paralleli:
        \[
        \mathbb{H} \cap \mathbb{H}' = \emptyset \quad \text{e} \quad \mathbb{H} \nparallel \mathbb{H}'
        \]
        \item Si dicono \textbf{strettamente sghembi} se non si intersecano e l'intersezione delle loro giaciture è solo il vettore nullo:
        \[
        \mathbb{H} \cap \mathbb{H}' = \emptyset \quad \text{e} \quad \vec{\mathbb{H}} \cap \vec{\mathbb{H}'} = \{\bar{0}\}
        \]
    \end{itemize}
}

\section{Classificazione e Rappresentazione Cartesiana}

Sia \((\vec{\mathcal{E}}, \mathcal{E}, \pi)\) uno spazio affine con \(\dim(\mathcal{E}) = n\).
Sia \(\mathbb{H} \subseteq \mathcal{E}\) un sottospazio affine di dimensione \(\dim(\mathbb{H}) = h\).
Usiamo la seguente nomenclatura classica:
\begin{itemize}
    \item Se \(h = 0\), \(\mathbb{H}\) è un \textbf{punto}.
    \item Se \(h = 1\), \(\mathbb{H}\) è una \textbf{retta}.
    \item Se \(h = 2\), \(\mathbb{H}\) è un \textbf{piano}.
    \item Se \(h = n-1\), \(\mathbb{H}\) è un \textbf{iperpiano}.
\end{itemize}

\Teorema{Rappresentazione Cartesiana di un Sottospazio Affine}{
    Sia \(\mathbb{H}\) un sottospazio affine di dimensione \(h\) in uno spazio \(\mathcal{E}\) di dimensione \(n\).
    Fissato un riferimento affine \(R = (O, \mathcal{B})\), esiste un sistema lineare \(\Sigma: AX = b\) in \(n\) incognite tale che l'insieme \(S\) delle sue soluzioni è costituito dai vettori delle coordinate in \(K\) dei punti di \(\mathbb{H}\).
    
    Inoltre, il sistema omogeneo associato \(AX = 0\) rappresenta la giacitura \(\vec{\mathbb{H}}\) e \(\operatorname{rango}(A) = n-h\).

    \Dimostrazione{
        Fissiamo un punto \(P \in \mathbb{H}\). Sappiamo che:
        \[ \mathbb{H} = (P, \vec{\mathbb{H}}) = \{Q \in \mathcal{E} \mid \vec{PQ} \in \vec{\mathbb{H}}\} \]
        
        Siano date le coordinate nel riferimento \(R\):
        \begin{itemize}
            \item Del punto fisso: \(P \equiv_R (a_1, \dots, a_n)\)
            \item Del punto generico: \(Q \equiv_R (y_1, \dots, y_n)\)
        \end{itemize}
        
        Seguiamo la catena di equivalenze:
        
        \[ Q \in \mathbb{H} \iff \vec{PQ} \in \vec{\mathbb{H}} \]
        
        Passando alle coordinate tramite l'isomorfismo \(\phi_{\mathcal{B}}\):
        \[ \iff \phi_{\mathcal{B}}(\vec{PQ}) \in \phi_{\mathcal{B}}(\vec{\mathbb{H}}) \subseteq K^n \]
        
        Per il teorema di rappresentazione dei sottospazi vettoriali, esiste un sistema omogeneo \(\Sigma_0: AX=0\) tale che il suo spazio delle soluzioni \(S_0\) coincide con le coordinate della giacitura (\(S_0 = \phi_{\mathcal{B}}(\vec{\mathbb{H}})\)), con \(\dim(S_0)=h\) e \(\operatorname{rango}(A)=n-h\).
        
        Quindi la condizione diventa:
        \[ \iff A \left( \phi_{\mathcal{B}}(\vec{PQ}) \right) = 0 \]
        
        Sapendo che le coordinate del vettore \(\vec{PQ}\) sono la differenza delle coordinate dei punti (\(Y - A_0\)):
        \[ \iff A \begin{pmatrix} y_1 - a_1 \\ \vdots \\ y_n - a_n \end{pmatrix} = 0 \]
        
        Sviluppando il prodotto matriciale:
        \[ \iff A \begin{pmatrix} y_1 \\ \vdots \\ y_n \end{pmatrix} - A \begin{pmatrix} a_1 \\ \vdots \\ a_n \end{pmatrix} = 0 \]
        
        Portando il termine noto a destra:
        \[ \iff A \begin{pmatrix} y_1 \\ \vdots \\ y_n \end{pmatrix} = A \begin{pmatrix} a_1 \\ \vdots \\ a_n \end{pmatrix} \]
        
        Ponendo \(b = A \cdot (a_1 \dots a_n)^t\) (che è un vettore costante), abbiamo dimostrato che:
        \[ (y_1, \dots, y_n) \text{ è soluzione di } AX = b \]
    }
}

\BoxGreen{Esempio}{Rappresentazione di una retta nel piano}{
    Sia \( (\vec{\mathcal{E}}, \mathcal{E}, \pi) \) uno spazio affine con \(\dim \mathcal{E} = 2\) e riferimento \(R = (O, \mathcal{B})\).
    Vogliamo rappresentare la retta \(r\) passante per i punti \(P(2,1)\) e \(P'(3, -1)\).

    \textbf{1. Calcolo della Giacitura}
    La retta è definita da un punto di passaggio (es. \(P\)) e dalla sua direzione (giacitura) \(\vec{r}\):
    \[ r = P + \vec{r}, \quad \text{dove } \vec{r} = L(\vec{PP'}) \]
    Calcoliamo le componenti del vettore direttore rispetto alla base \(\mathcal{B}\):
    \[ \phi_{\mathcal{B}}(\vec{PP'}) = (3-2, -1-1) = (1, -2) \]
    Quindi la giacitura è generata dal vettore \( v = (1, -2) \).

    \textbf{2. Metodo Cartesiano (tramite Rango/Determinante)}
    Un punto generico \(Q(x_1, x_2)\) appartiene a \(r\) se e solo se il vettore \(\vec{PQ}\) appartiene alla giacitura \(\vec{r}\) (cioè è parallelo a \(v\)).
    \[ \vec{PQ} \equiv (x_1 - 2, x_2 - 1) \]
    La condizione è che i vettori \(\vec{PQ}\) e \(v\) siano linearmente dipendenti. Costruiamo la matrice delle coordinate:
    \[
    A = \begin{pmatrix} 1 & x_1 - 2 \\ -2 & x_2 - 1 \end{pmatrix}
    \]
    Imponiamo che il rango sia 1 (cioè determinante nullo):
    \[ \det(A) = 1(x_2 - 1) - (-2)(x_1 - 2) = 0 \]
    \[ x_2 - 1 + 2x_1 - 4 = 0 \]
    \[ 2x_1 + x_2 - 5 = 0 \quad (\text{Equazione Cartesiana}) \]

    \textbf{3. Metodo Parametrico}
    Alternativamente, \(Q \in r\) se esiste un parametro \(t \in \mathbb{R}\) tale che \(\vec{PQ} = t \cdot v\).
    \[ (x_1 - 2, x_2 - 1) = t(1, -2) \]
    Uguagliando componente per componente otteniamo il sistema parametrico:
    \[
    \begin{cases}
    x_1 - 2 = t \\
    x_2 - 1 = -2t
    \end{cases}
    \implies
    \begin{cases}
    x_1 = 2 + t \\
    x_2 = 1 - 2t
    \end{cases}
    \]
    \textit{Passaggio alla forma cartesiana:} Ricaviamo \(t\) dalla prima equazione e sostituiamo nella seconda:
    \[ t = x_1 - 2 \implies x_2 = 1 - 2(x_1 - 2) \]
    \[ x_2 = 1 - 2x_1 + 4 \implies 2x_1 + x_2 - 5 = 0 \]
}

\BoxGreen{Esempio}{Rappresentazione di una retta nello spazio}{
    Sia \( (\vec{\mathcal{E}}, \mathcal{E}, \pi) \) uno spazio affine con \(\dim \mathcal{E} = 3\) e riferimento \(R = (O, \mathcal{B})\).
    Vogliamo rappresentare la retta \(r\) passante per \(P(1, 2, -2)\) e \(P'(3, 3, 1)\).

    \textbf{1. Calcolo della Giacitura (Direzione)}
    La giacitura della retta è generata dal vettore differenza tra i due punti:
    \[ \vec{r} = L(\vec{PP'}) \]
    Calcoliamo le componenti:
    \[ \vec{v} = \vec{PP'} = (3-1, \ 3-2, \ 1-(-2)) = (2, 1, 3) \]

    \textbf{2. Rappresentazione Parametrica}
    Un generico punto \(Q(x_1, x_2, x_3)\) appartiene alla retta se e solo se il vettore \(\vec{PQ}\) è multiplo di \(\vec{v}\).
    \[ \vec{PQ} = t \cdot \vec{v} \implies (x_1-1, x_2-2, x_3+2) = t(2, 1, 3) \]
    Uguagliando le componenti, otteniamo il sistema parametrico:
    \[
    \begin{cases}
    x_1 = 1 + 2t \\
    x_2 = 2 + t \\
    x_3 = -2 + 3t
    \end{cases}
    \quad \text{con } t \in \mathbb{R}
    \]

    \textbf{3. Rappresentazione Cartesiana (Eliminazione del parametro)}
    Per trovare le equazioni cartesiane, ricaviamo \(t\) da una delle equazioni (ad esempio la seconda, che è la più semplice) e sostituiamo nelle altre due.
    
    Dalla seconda equazione: \( t = x_2 - 2 \).
    
    Sostituiamo nella prima:
    \[ x_1 = 1 + 2(x_2 - 2) \implies x_1 = 1 + 2x_2 - 4 \implies x_1 - 2x_2 + 3 = 0 \]
    
    Sostituiamo nella terza:
    \[ x_3 = -2 + 3(x_2 - 2) \implies x_3 = -2 + 3x_2 - 6 \implies -3x_2 + x_3 + 8 = 0 \]
    
    Il sistema cartesiano \(\Sigma\) che descrive la retta è l'intersezione di due piani:
    \[
    \Sigma : \begin{cases}
    x_1 - 2x_2 + 3 = 0 \\
    -3x_2 + x_3 + 8 = 0
    \end{cases}
    \]

    \textbf{4. Verifica (Giacitura)}
    Il sistema omogeneo associato \(\Sigma_0\) descrive la giacitura \(\vec{r}\):
    \[
    \Sigma_0 : \begin{cases}
    x_1 - 2x_2 = 0 \\
    -3x_2 + x_3 = 0
    \end{cases}
    \]
    Possiamo verificare se il vettore direttore \((2, 1, 3)\) soddisfa questo sistema:
    \begin{itemize}
        \item \(2 - 2(1) = 0\) (Vero)
        \item \(-3(1) + 3 = 0\) (Vero)
    \end{itemize}
    I conti tornano.
}

\BoxGreen{Esempio}{Rappresentazione Cartesiana di un Piano}{
    Sia \( (\vec{\mathcal{E}}, \mathcal{E}, \pi) \) uno spazio affine con \(\dim \mathcal{E} = 3\) e riferimento \(R=(O, \mathcal{B})\).
    Vogliamo rappresentare il piano \(\mathbb{H}\) passante per i tre punti:
    \[ P_0=(0,-2,1), \quad P_1=(1,1,-1), \quad P_2=(2,1,0) \]

    \textbf{1. Calcolo della Giacitura}
    Il piano è individuato da un punto (es. \(P_0\)) e dalla giacitura generata dai vettori che collegano i punti:
    \[ \vec{\mathbb{H}} = L(\vec{P_0 P_1}, \vec{P_0 P_2}) \]
    Calcoliamo le componenti dei vettori direttori:
    \begin{itemize}
        \item \(\vec{P_0 P_1} = (1-0, 1-(-2), -1-1) = (1, 3, -2)\)
        \item \(\vec{P_0 P_2} = (2-0, 1-(-2), 0-1) = (2, 3, -1)\)
    \end{itemize}
    Verifica di indipendenza: i due vettori non sono proporzionali, quindi sono linearmente indipendenti. Ciò implica che \(P_0, P_1, P_2\) sono \textbf{affinemente indipendenti} (non allineati) e generano effettivamente un piano (dimensione 2).

    \textbf{2. Costruzione della Matrice (Rouché-Capelli)}
    Un punto generico \(Q(x_1, x_2, x_3)\) appartiene al piano \(\mathbb{H}\) se e solo se il vettore \(\vec{P_0 Q}\) appartiene alla giacitura \(\vec{\mathbb{H}}\) (cioè è combinazione lineare di \(\vec{P_0 P_1}\) e \(\vec{P_0 P_2}\)).
    
    Calcoliamo \(\vec{P_0 Q} = (x_1 - 0, \ x_2 - (-2), \ x_3 - 1) = (x_1, x_2+2, x_3-1)\).
    
    La condizione è:
    \[
    Q \in \mathbb{H} \iff \operatorname{rango} \begin{pmatrix} 
    1 & 2 & x_1 \\ 
    3 & 3 & x_2+2 \\ 
    -2 & -1 & x_3-1 
    \end{pmatrix} = 2
    \]

    \textbf{3. Riduzione a Gradini}
    Riduciamo la matrice per trovare l'equazione (imponendo che l'ultima riga si annulli).
    \[
    \begin{pmatrix} 1 & 2 & x_1 \\ 3 & 3 & x_2+2 \\ -2 & -1 & x_3-1 \end{pmatrix}
    \xrightarrow[R_3 \to R_3 + 2R_1]{R_2 \to R_2 - 3R_1}
    \begin{pmatrix} 1 & 2 & x_1 \\ 0 & -3 & (x_2+2) - 3x_1 \\ 0 & 3 & (x_3-1) + 2x_1 \end{pmatrix}
    \]
    
    Procediamo con l'ultimo passaggio per eliminare il pivot nella seconda colonna:
    \[
    \xrightarrow{R_3 \to R_3 + R_2}
    \begin{pmatrix} 
    1 & 2 & x_1 \\ 
    0 & -3 & -3x_1 + x_2 + 2 \\ 
    0 & 0 & \mathbf{(x_3 - 1 + 2x_1) + (-3x_1 + x_2 + 2)} 
    \end{pmatrix}
    \]

    \textbf{4. Equazione Finale}
    Affinché il rango sia 2, l'elemento nell'ultima riga e ultima colonna deve essere nullo:
    \[
    x_3 - 1 + 2x_1 - 3x_1 + x_2 + 2 = 0
    \]
    \[
    -x_1 + x_2 + x_3 + 1 = 0
    \]
    
    Questa è la rappresentazione cartesiana del piano \(\mathbb{H}\).
}

\Proposizione{Sistemi Lineari e Sottospazi Affini}{
    Sia \( \Sigma: AX = b \) un sistema lineare compatibile su un campo \( K \) in \( n \) incognite.
    Sia \( \rho = \operatorname{rango}(A) \).
    
    Allora l'insieme delle soluzioni \( S \) del sistema rappresenta le coordinate di un \textbf{sottospazio affine} \( \mathbb{H} \) di uno spazio affine \( (\vec{\mathcal{E}}, \mathcal{E}, \pi) \) di dimensione \( n \), rispetto a un riferimento fissato \( R = (O, \mathcal{B}) \).
    
    La dimensione del sottospazio è data dalla codimensione del rango:
    \[ \dim(\mathbb{H}) = n - \rho \]

    \Dimostrazione{
        Dal teorema di struttura dei sistemi lineari, sappiamo che la soluzione generale si scrive come somma di una soluzione particolare e della soluzione del sistema omogeneo associato.
        
        \begin{enumerate}
            \item \textbf{Punto di passaggio (Soluzione Particolare):}
            Poiché il sistema è compatibile, esiste almeno una soluzione \( (a_1, \dots, a_n) \).
            Sia \( P \) il punto dello spazio affine che ha queste coordinate nel riferimento \( R \):
            \[ P \equiv_R (a_1, \dots, a_n) \]
            
            \item \textbf{Giacitura (Sistema Omogeneo):}
            Consideriamo il sistema omogeneo associato \( \Sigma_0: AX = 0 \).
            L'insieme delle sue soluzioni \( S_0 \) è un sottospazio vettoriale di \( K^n \).
            Definiamo la giacitura \( \vec{\mathbb{H}} \) come il sottospazio di \( \vec{\mathcal{E}} \) che corrisponde a \( S_0 \) tramite l'isomorfismo inverso delle coordinate:
            \[ \vec{\mathbb{H}} = \phi_{\mathcal{B}}^{-1}(S_0) \]
            Per il teorema del rango (applicato ai sistemi omogenei), la dimensione dello spazio delle soluzioni è il numero di incognite meno il rango della matrice:
            \[ \dim(\vec{\mathbb{H}}) = \dim(S_0) = n - \rho \]
        \end{enumerate}
        
        Pertanto, l'insieme delle soluzioni di \( \Sigma \) descrive il sottospazio affine passante per \( P \) con giacitura \( \vec{\mathbb{H}} \):
        \[ \mathbb{H} = (P, \vec{\mathbb{H}}) = P + \vec{\mathbb{H}} \]
    }
}

\Esempio{Calcolo di un Iperpiano in \(\mathbb{R}^4\)}{
    Sia \( \dim(\mathcal{E}) = 4 \). Consideriamo il sottospazio affine \( \mathbb{H} \) definito dall'equazione cartesiana:
    \[ \Sigma: 3x_1 - x_2 + 4x_3 + x_4 = 2 \]
    
    \textbf{1. Analisi del Rango e Dimensione}
    La matrice dei coefficienti è \( A = (3, -1, 4, 1) \).
    Il rango è \(\rho = 1\) (c'è almeno un elemento non nullo).
    La dimensione di \( \mathbb{H} \) sarà:
    \[ \dim(\mathbb{H}) = n - \rho = 4 - 1 = 3 \]
    (Si tratta quindi di un \textit{iperpiano} affine).

    \textbf{2. Calcolo della Giacitura (\(\vec{\mathbb{H}}\))}
    Risolviamo il sistema omogeneo associato \( \Sigma_0: 3x_1 - x_2 + 4x_3 + x_4 = 0 \).
    Possiamo esplicitare \( x_2 \) (che ha coefficiente -1, comodo per i calcoli):
    \[ x_2 = 3x_1 + 4x_3 + x_4 \]
    Le variabili libere sono \( x_1, x_3, x_4 \). Assegniamo loro i parametri per trovare una base di \( S_0 \):
    \begin{itemize}
        \item \( x_1=1, x_3=0, x_4=0 \implies x_2 = 3 \cdot 1 = 3 \to v_1 = (1, 3, 0, 0) \)
        \item \( x_1=0, x_3=1, x_4=0 \implies x_2 = 4 \cdot 1 = 4 \to v_2 = (0, 4, 1, 0) \)
        \item \( x_1=0, x_3=0, x_4=1 \implies x_2 = 1 \cdot 1 = 1 \to v_3 = (0, 1, 0, 1) \)
    \end{itemize}
    La giacitura è \( \vec{\mathbb{H}} = L(v_1, v_2, v_3) \).

    \textbf{3. Punto di Passaggio}
    Cerchiamo una soluzione particolare del sistema completo (non omogeneo).
    Poniamo le variabili libere a 0 (\(x_1=0, x_3=0, x_4=0\)):
    \[ -x_2 = 2 \implies x_2 = -2 \]
    Un punto di passaggio è \( P \equiv (0, -2, 0, 0) \).
    
    \textbf{Conclusione:}
    Il sottospazio è \( \mathbb{H} = P + L(v_1, v_2, v_3) \).
}